\documentclass{situationist-preprint}

\title{Self-Locating Belief and the Indexical \enquote{Now}}
\msid{EPI-18-c}
\msfamily{EPI-18}
\msdiscipline{Epistemology}
\mstype{Consequence}
\mskeywords{Sleeping Beauty; self-locating belief; indexicals; now; B-theory; centred worlds; tensed facts}

\begin{document}
\maketitle

\begin{abstract}
In the Sleeping Beauty problem, Beauty wonders whether it is now Monday, and on being told that it is, her credence in heads moves from one third to one half. It may seem that this requires tensed facts: that \enquote{it is now Monday} reports a fact about which time is present, which Beauty lacks on waking and learns later. We show that the thirder solution requires nothing of the kind. On a tenseless view of time, on which \enquote{now} is an indexical referring to the time of the thought, Beauty's uncertainty is uncertainty about where she is located in a world whose history she knows apart from the coin. Her credences are defined on centred worlds, her learning is conditionalization on self-locating information, and her loss of information between Sunday and waking is a change of centre. The thirder answer and the claim that the residual disagreement concerns the content of self-locating belief both follow, and the A-theorist's tensed facts do no work in the argument.
\end{abstract}

\section{The question}

EPI-18-a argues that, under the standard protocol in which her awakenings are subjectively indistinguishable, Beauty's credence in heads on waking is one third, and that the residual disagreement with the halfer concerns the content of self-locating belief. The argument is stated in terms of centred worlds \citep{Lewis1979}, and its key premise, the restricted Principal Principle (Principle~P1 there), is applied to Beauty's credence in heads \emph{conditional on its being Monday} (Lemma~3.1). On being told \enquote{it is now Monday}, Beauty conditionalizes and her credence in heads rises from one third to one half (Section~8 there).

What does she learn? An A-theorist says she learns a tensed fact: that Monday is present. On the A-theory there is an objective present, and \enquote{it is now Monday} is true just when Monday has the property of presentness. If that were so, Beauty's problem would be a problem about the tensed facts, and the thirder solution would depend on the metaphysics of tense.

We take the B-theory as given. Time is constituted by the B-series of earlier and later (TSM-01); the present has no metaphysical privilege, being the time at which a thought or utterance occurs (TSM-15); and \enquote{now} is an indexical, so that the question why it is now now has a trivial answer (TSM-18). Our aim is to derive what the thirder solution then says about the content of Beauty's beliefs, and to show that nothing in the solution is lost.

\section{Tenseless content}

\begin{definition}[Centres and centred worlds]\label{d:centre}
A \emph{centre} in a world $w$ is a pair of a subject and a time at which she is thinking. A \emph{centred world} is a world together with a centre in it. A \emph{self-locating proposition} is a set of centred worlds that may contain some but not all of the centres of a world.
\end{definition}

\begin{assumption}[Indexical \enquote{now}]\label{a:now}
A token of \enquote{now} in a thought refers to the time at which the thought occurs. A thought token of the form \enquote{it is now $D$} is true if and only if the token occurs on day $D$.
\end{assumption}

Assumption~\ref{a:now} is the token-reflexive account of tensed thought and talk that the B-theorist offers \citep{Mellor1998}. It gives the truth conditions of each token tenselessly, by the date of the token. The belief that it is now Monday is then not the belief that a tensed fact obtains. It is the self-ascription of a property: that of thinking on a Monday \citep{Lewis1979}.

\begin{proposition}[No tensed ignorance]\label{p:notensed}
On waking, Beauty is ignorant of no tenseless fact other than the outcome of the coin, and of no tensed fact at all. Her remaining uncertainty is uncertainty about which centre she occupies.
\end{proposition}

\begin{proof}[Argument]
Beauty knows the protocol, and so knows, for each outcome of the coin, the complete B-series history of the experiment: on which days she is woken, and when her memory is erased. On the B-theory there are no further facts about which time is present (TSM-15). Her credence on waking is therefore distributed over the three centred worlds $H_1$, $T_1$ and $T_2$ of EPI-18-a: two uncentred possibilities, one of which contains two centres. What she lacks, beyond the outcome of the coin, is knowledge of which of these centres is hers.
\end{proof}

\begin{corollary}[What \enquote{it is Monday} adds]\label{c:monday}
When Beauty is told that it is now Monday, she learns the self-locating proposition $\{H_1, T_1\}$. The uncentred proposition that she is awake on Monday was already certain, and she learns no tensed fact.
\end{corollary}

\begin{proof}
By Assumption~\ref{a:now}, the information that it is now Monday is true at a centre just when the centre's time falls on Monday; these are $H_1$ and $T_1$. The uncentred proposition that Beauty is woken on Monday is true at both the heads and the tails world, so its probability was one. By Proposition~\ref{p:notensed} there is no tensed fact to learn.
\end{proof}

Corollary~\ref{c:monday} is a consequence with some interest of its own. Beauty's credence in heads changes from one third to one half on information that is, considered as uncentred, certain in advance. The change cannot be explained by uncentred learning, and on the B-theory it cannot be explained by tensed learning either. It is explained by self-locating learning: of the three centres she might occupy, she learns that she occupies one of two, and conditionalization on this proposition gives
\[
P(H_1 \mid \{H_1, T_1\}) = \frac{P(H_1)}{P(H_1) + P(T_1)} = \frac{1}{2},
\]
as in EPI-18-a. Self-locating information bears on uncentred hypotheses whenever the hypotheses differ in how many centres satisfy it \citep{Elga2004}, and nothing in this requires that some time be objectively present.

\section{Change without passage}

The A-theorist may press the following point. Beauty's credence that it is now Monday is one third on waking on Monday and still one third on waking on Tuesday (given tails), but the proposition it is about is true on the first occasion and false on the second. If the object of her belief is the same in both, it must be a proposition whose truth value changes with time, and such a proposition is tensed \citep{Prior1959}.

\begin{proposition}[Change of centre]\label{p:shift}
The same self-locating proposition can be true at one of Beauty's centres and false at another without any change in what the world is like. Her credence function on Tuesday differs from her credence function on Monday only in being had at a different centre.
\end{proposition}

\begin{proof}[Argument]
By Definition~\ref{d:centre}, a self-locating proposition is true or false at a centred world, not at a world. $\{H_1, T_1\}$ is true at $T_1$ and false at $T_2$; both centres belong to the tails world, which is the same tenseless world throughout. Beauty's credence function is a function on centred worlds, and on each waking it is the same function, since by Assumption~2.1 of EPI-18-a her evidence is the same. What differs between Monday and Tuesday is which centre she occupies, which is a tenseless fact about the location of a thought in the B-series.
\end{proof}

The proposition that seems to change its truth value is a property, true of some centres and false of others. This is the B-theorist's general account of the phenomena that Prior took to require tensed facts: the relief that a headache is over is directed at a self-locating content and is had at a time after the headache, and the B-theorist explains it by those two facts together \citep{Mellor1998}. The impression that something changes in the world as Beauty passes from Monday to Tuesday is an impression of her own changing location, which is what TSM-02 says of the impression of passage in general.

The same account applies to the loss of information that EPI-18-a uses in its treatment of diachronic coherence (Section~8 there). On Sunday Beauty knew that it was Sunday; on waking she does not know the day. On the A-theory this would be the loss of a tensed fact. On the present account it is the loss of self-locating knowledge: her Sunday centre was located with certainty, her waking centre is not. That loss is not a case for conditionalization \citep{Arntzenius2003}, and it requires no more than that her later centre lacks some evidence that her earlier one had.

\section{Objections}

\begin{objection}
Perry's essential indexical shows that indexical beliefs explain action in a way that no non-indexical belief can. If Beauty's belief that it is now Monday were tenseless, it could not explain why she acts as she does on Monday.
\end{objection}

\begin{reply}
Perry's argument shows that indexical beliefs cannot be replaced by beliefs in uncentred propositions \citep{Perry1979}. It does not show that they are beliefs in tensed facts. A self-locating belief is essentially indexical in Perry's sense: it is the self-ascription of a property, and it explains action because the subject who ascribes it ascribes it to herself, at the time of acting. The B-theory supplies such contents, and Beauty's self-locating credences are of that kind.
\end{reply}

\begin{objection}
If there are no tensed facts, Beauty's question \enquote{is it now Monday?} has no factual answer, and her credence in it is not a credence about anything.
\end{objection}

\begin{reply}
The question has a factual answer at each centre at which it is asked: on the token-reflexive account, it is true at Monday centres and false at Tuesday centres. Her credence measures her uncertainty about which centre she occupies, and the fact that settles it is the tenseless date of her thought. Nothing more is needed for credence to be about something.
\end{reply}

\begin{objection}
The halfer may grant all this and conclude that self-locating credence, being about centres rather than worlds, should not bear on her credence in heads at all.
\end{objection}

\begin{reply}
That is the residual disagreement about the content of self-locating belief, which EPI-18-a identifies (Section~10 there), and which EPI-18-b examines from the halfer's side. The B-theory does not settle it, but it removes one source of confusion from it. The halfer cannot claim that Beauty's credence in heads is about tenseless worlds while her credence about the day is about tensed facts; on the B-theory both are credences over centred worlds, and the restricted Principal Principle, applied to her credence in heads conditional on the self-locating proposition $\{H_1, T_1\}$, yields the thirder answer by Theorem~3.3 of EPI-18-a.
\end{reply}

\section{Conclusion}

On a tenseless view of time, \enquote{now} is an indexical and \enquote{it is now Monday} expresses self-locating information. Beauty's uncertainty on waking is uncertainty about her location in a B-series she otherwise knows; what she learns when told the day is a self-locating proposition; and what she lost between Sunday and waking is self-locating knowledge. The thirder solution of EPI-18-a needs no tensed facts, and the disagreement that remains concerns the content of self-locating belief, not the metaphysics of the present.

\begin{thebibliography}{99}\small
\bibitem[Arntzenius(2003)]{Arntzenius2003} Arntzenius, F. (2003). Some problems for conditionalization and reflection. \emph{The Journal of Philosophy}, 100(7), 356--370.
\bibitem[Elga(2004)]{Elga2004} Elga, A. (2004). Defeating Dr. Evil with self-locating belief. \emph{Philosophy and Phenomenological Research}, 69(2), 383--396.
\bibitem[Lewis(1979)]{Lewis1979} Lewis, D. (1979). Attitudes de dicto and de se. \emph{The Philosophical Review}, 88(4), 513--543.
\bibitem[Mellor(1998)]{Mellor1998} Mellor, D.~H. (1998). \emph{Real Time II}. Routledge.
\bibitem[Perry(1979)]{Perry1979} Perry, J. (1979). The problem of the essential indexical. \emph{No\^us}, 13(1), 3--21.
\bibitem[Prior(1959)]{Prior1959} Prior, A.~N. (1959). Thank goodness that's over. \emph{Philosophy}, 34(128), 12--17.
\end{thebibliography}

\end{document}
